\begin{exo}{}
	On considère les vecteurs suivants dans un repère :

	\begin{center}
		\begin{tikzpicture}[>=Stealth]
			\draw[gray!30, thick, step=1] (0.5,0.5) grid (8.5,7.5);
			\coordinate (A) at (1,3);
			\coordinate (B) at (5,6);
			\coordinate (C) at (8,5);
			\coordinate (D) at (4,2);
			\coordinate (E) at (2,5);
			\coordinate (F) at (6,5);
			\coordinate (G) at (3,2);
			\coordinate (H) at (4,7);
			\coordinate (I) at (3,7);
			\coordinate (J) at (6,7);
			\coordinate (K) at (4,1);
			\coordinate (L) at (7,1);
			\coordinate (M) at (4,4);
			\coordinate (N) at (4,3);
			\draw (A) node[below] {$A$} node[gray] {$\bullet$};
			\draw (B) node[above] {$B$} node[gray] {$\bullet$};
			\draw (C) node[above] {$C$} node[gray] {$\bullet$};
			\draw (D) node[below] {$D$} node[gray] {$\bullet$};
			\draw (E) node[left] {$E$} node[gray] {$\bullet$};
			\draw (F) node[right] {$F$} node[gray] {$\bullet$};
			\draw (G) node[below] {$G$} node[gray] {$\bullet$};
			\draw (H) node[above] {$H$} node[gray] {$\bullet$};
			\draw (I) node[left] {$I$} node[gray] {$\bullet$};
			\draw (J) node[right] {$J$} node[gray] {$\bullet$};
			\draw (K) node[left] {$K$} node[gray] {$\bullet$};
			\draw (L) node[right] {$L$} node[gray] {$\bullet$};
			\draw (M) node[above] {$M$} node[gray] {$\bullet$};
			\draw (N) node[below] {$N$} node[gray] {$\bullet$};
			\draw[very thick, ->] (A) -- (B);
			\draw[very thick, ->] (C) -- (D);
			\draw[very thick, ->] (E) -- (F);
			\draw[very thick, ->] (E) -- (G);
			\draw[very thick, ->] (H) -- (I);
			\draw[very thick, ->] (H) -- (J);
			\draw[very thick, ->] (K) -- (L);
			\draw[very thick, ->] (M) -- (N);
		\end{tikzpicture}
	\end{center}

	\begin{enumerate}
		\item \textbf{Donner trois phrases} en vous appuyant sur les termes suivants : \emph{extrémité}, \emph{origine}, \emph{norme}.
		\item \textbf{Indiquer} les vecteurs de même direction, puis ceux de même sens, puis de sens opposé et enfin de même norme.
	\end{enumerate}
\end{exo}
